\chapter{Orthogonalization: Arnoldi, Lanczos, and Gram--Schmidt}
\label{ch:orthog}

\chapterorient{Every Krylov solver in this Part rests on one ability: to take a
growing pile of nearly parallel vectors and straighten them into an orthonormal
basis. GMRES (\cref{ch:gmres}) called this step \code{orthogonalize} and moved
on; this is the careful version. We build Gram--Schmidt from its geometry, then
meet its two faces---classical and modified---and discover that the difference
between them is not arithmetic but \emph{order}: modified Gram--Schmidt is a
genuine sequential dependency, a concrete specimen of the obstruction of
\cref{ch:rotations}. We watch both lose orthogonality in finite precision, learn
the ``twice is enough'' repair, and then assemble the orthogonalizer into the two
recurrences the rest of the book runs on: Arnoldi, which feeds GMRES, and
Lanczos, which---for symmetric operators---collapses to a miraculous three-term
recurrence and then, just as miraculously, falls apart numerically. The thread
running through all of it is a single \emph{contract}: ``the basis is
orthonormal, to a tolerance.'' Everything downstream depends on that contract
holding.}

% local: lowercase bold-a for the worked-example input vectors (not in the shared preamble)
\providecommand{\va}{\mathbf{a}}

\section{The goal: an orthonormal basis for a Krylov subspace}
\label{sec:goal}

A Krylov subspace solver does not work with the operator $\mat{A}$ directly. It
works inside the \emph{Krylov subspace}
\[
  \mathcal{K}_m(\mat{A}, \vv)
    \;=\; \operatorname{span}\bigl\{\vv,\;\mat{A}\vv,\;\mat{A}^2\vv,\;\dots,\;
      \mat{A}^{m-1}\vv\bigr\},
\]
the span of the vector $\vv$ and its repeated images under $\mat{A}$. The reason
this subspace is useful is the subject of \cref{ch:gmres} and \cref{ch:cg}; here
we take it as given and ask a narrower question. To compute \emph{inside}
$\mathcal{K}_m$---to solve a least-squares problem there, to project the operator
onto it---we need a \emph{basis}: a set of $m$ vectors that span it and that we
can compute with. The obvious basis is the one in the definition, the
\emph{power basis} $\{\vv, \mat{A}\vv, \mat{A}^2\vv, \dots\}$. It is also,
catastrophically, the worst possible basis to compute with.

\begin{figure}[t]
\centering
\begin{tikzpicture}[font=\footnotesize, scale=1.0]
  % ===== LEFT: the power basis collapses toward a line =====
  \begin{scope}
    \node[font=\small\bfseries, text=simRust] at (1.5,2.7) {power basis: nearly parallel};
    \draw[->,simGray] (-0.2,0)--(3.1,0);
    \draw[->,simGray] (0,-0.2)--(0,2.4);
    % v, Av, A2v all crowding one direction
    \draw[flow,simRust] (0,0)--(2.6,1.5) node[right,font=\scriptsize]{$\vv$};
    \draw[flow,simRust] (0,0)--(2.75,1.25) node[right,font=\scriptsize]{$\mat{A}\vv$};
    \draw[flow,simRust] (0,0)--(2.85,1.05) node[right,font=\scriptsize]{$\mat{A}^2\vv$};
    \node[font=\scriptsize\itshape, text=simGray, align=center] at (1.5,-0.7)
      {they crowd one direction:\\the dominant eigenvector};
  \end{scope}
  \draw[simGray!45, thick] (3.7,-1.0) -- (3.7,3.0);
  % ===== RIGHT: the orthonormal basis spreads out =====
  \begin{scope}[xshift=5.0cm]
    \node[font=\small\bfseries, text=simGreen] at (1.3,2.7) {orthonormal basis: spread out};
    \draw[->,simGray] (-0.2,0)--(2.6,0);
    \draw[->,simGray] (0,-0.2)--(0,2.4);
    \draw[flow,simGreen] (0,0)--(2.0,0) node[right,font=\scriptsize]{$\vq_1$};
    \draw[flow,simGreen] (0,0)--(0,2.0) node[above,font=\scriptsize]{$\vq_2$};
    \draw (1.7,0) -- (1.7,0.3) -- (2.0,0.3);
    \node[font=\scriptsize] at (1.9,0.5) {$90^\circ$};
    \node[font=\scriptsize\itshape, text=simGray, align=center] at (1.3,-0.7)
      {mutually orthogonal,\\unit length: well-conditioned};
  \end{scope}
\end{tikzpicture}
\caption[Power basis versus orthonormal basis]{Why the Krylov subspace needs an
orthonormal basis. \emph{Left:} the natural power basis
$\{\vv, \mat{A}\vv, \mat{A}^2\vv, \dots\}$ is a disaster to compute with---each
application of $\mat{A}$ pulls the vector further toward the dominant eigenvector,
so successive members crowd into nearly the same direction. A matrix whose columns
are these vectors is severely ill-conditioned, and in finite precision the later
members are numerically indistinguishable from the earlier ones. \emph{Right:} an
\emph{orthonormal} basis $\{\vq_1, \vq_2, \dots\}$ spanning the same subspace
spreads the directions out to mutual right angles and unit length. It is the basis
every Krylov method actually uses, and producing it from the power basis is the
business of this chapter.}
\label{fig:powerbasis}
\end{figure}

The trouble is visible in \cref{fig:powerbasis}. Each multiplication by $\mat{A}$
amplifies the component of the vector along the \emph{dominant} eigenvector---the
one with the largest eigenvalue in magnitude---more than any other. So $\mat{A}\vv$
leans toward that eigenvector, $\mat{A}^2\vv$ leans further, and within a handful
of steps the power-basis vectors are all crowding into nearly the same direction.
They still, in exact arithmetic, span the full subspace; but the matrix whose
columns they are is wildly ill-conditioned, and in floating point the later columns
carry almost no new information. This is the same effect that makes the
\emph{power method} converge to the dominant eigenvector---a feature there, a bug
here.

The cure is to build, as we go, an \emph{orthonormal} basis for the same subspace:
a set of vectors $\vq_1, \vq_2, \dots, \vq_m$ that are mutually orthogonal and of
unit length,
\[
  \inner{\vq_i}{\vq_j} = \delta_{ij}
    \;=\;
    \begin{cases} 1 & i = j,\\ 0 & i \ne j, \end{cases}
\]
and that span $\mathcal{K}_m$ exactly as the power basis does. An orthonormal basis
is perfectly conditioned---the matrix $\mat{Q}$ with these columns satisfies
$\mat{Q}^{\!\top}\mat{Q} = \mat{I}$, the best-conditioned matrix there is---so every
later computation inside $\mathcal{K}_m$ is numerically clean. The procedure that
manufactures such a basis, one vector at a time, from the power basis is
\emph{Gram--Schmidt orthogonalization}, and it is the foundation everything in this
chapter is built on.

\begin{keyidea}
A Krylov solver works inside the subspace
$\mathcal{K}_m(\mat{A},\vv) = \operatorname{span}\{\vv, \mat{A}\vv, \dots,
\mat{A}^{m-1}\vv\}$, but never in the natural power basis, which collapses toward
the dominant eigenvector and is hopelessly ill-conditioned. It works instead in an
\emph{orthonormal} basis $\vq_1,\dots,\vq_m$ for the same subspace---mutually
orthogonal, unit length, $\inner{\vq_i}{\vq_j} = \delta_{ij}$. Building that basis
incrementally, as each new Krylov vector arrives, is the job of Gram--Schmidt, and
the quality of the basis it produces is what every downstream method silently
relies on.
\end{keyidea}

\section{Gram--Schmidt: subtract the projections, then normalize}
\label{sec:gschmidt}

Gram--Schmidt is one idea, applied repeatedly. Suppose we have already built an
orthonormal set $\vq_1, \dots, \vq_{m}$ and a new vector $\vw$ arrives (in the
Krylov setting, $\vw = \mat{A}\vq_m$, the next power-basis member). We want to turn
$\vw$ into the next basis vector $\vq_{m+1}$. The new vector almost certainly
overlaps the directions we already have; the orthonormal basis must point somewhere
\emph{new}. So we do two things, in order:

\begin{enumerate}
  \item \textbf{Subtract off the part of $\vw$ that lies in the span we already
  have.} The component of $\vw$ along a unit vector $\vq_j$ is the scalar
  projection $\inner{\vw}{\vq_j}$, and the vector projection is
  $\inner{\vw}{\vq_j}\,\vq_j$. Subtracting all of these leaves the
  \emph{orthogonal residual}
  \[
    \vw' \;=\; \vw \;-\; \sum_{j=1}^{m} \inner{\vw}{\vq_j}\,\vq_j,
  \]
  the part of $\vw$ that points outside $\operatorname{span}\{\vq_1,\dots,\vq_m\}$.
  \item \textbf{Normalize.} The residual $\vw'$ has the right \emph{direction} but
  not unit length, so we scale it:
  \[
    \vq_{m+1} \;=\; \frac{\vw'}{\norm{\vw'}}.
  \]
\end{enumerate}

The first step is the heart of the matter; the second is bookkeeping.
\Cref{fig:gschmidt} draws the geometry in three dimensions, where it can be seen
all at once. The vector $\vw$ casts a ``shadow'' onto the plane spanned by
$\vq_1, \vq_2$; subtracting that shadow leaves exactly the vertical component, which
points out of the plane in a brand-new direction.

\begin{figure}[t]
\centering
\begin{tikzpicture}[font=\footnotesize, scale=1.05, >={Stealth[length=2.4mm]}]
  % the q1-q2 plane drawn as a parallelogram
  \fill[simBgGreen, opacity=0.7] (-0.3,-0.4) -- (3.0,-0.1) -- (3.6,0.9) -- (0.3,0.6) -- cycle;
  \node[font=\scriptsize\itshape, text=simGreen] at (3.1,1.2) {$\operatorname{span}\{\vq_1,\vq_2\}$};
  % basis vectors q1, q2 in the plane (origin at O)
  \coordinate (O) at (0.6,0.1);
  \draw[flow,simGreen] (O) -- ++(2.0,0.15) node[below right,font=\scriptsize]{$\vq_1$};
  \draw[flow,simGreen] (O) -- ++(0.55,0.45) node[above left,font=\scriptsize]{$\vq_2$};
  % w sticking up out of the plane
  \coordinate (Wtip) at (2.3,2.5);
  \draw[flow,simRust,very thick] (O) -- (Wtip) node[above,font=\scriptsize]{$\vw$};
  % the shadow (projection onto the plane)
  \coordinate (Ptip) at (2.55,1.0);
  \draw[flow,simBlue] (O) -- (Ptip)
    node[pos=0.6,below right,font=\scriptsize,text=simBlue]
      {$\displaystyle\sum_j\inner{\vw}{\vq_j}\vq_j$};
  % the orthogonal residual w' (vertical), drawn from Ptip up to Wtip
  \draw[flow,simViolet,very thick] (Ptip) -- (Wtip)
    node[pos=0.5,right,font=\scriptsize,text=simViolet]{$\vw'$ (residual)};
  % dashed drop line emphasising the shadow
  \draw[densely dashed,simGray] (Wtip) -- (Ptip);
  % right-angle mark at Ptip between plane and residual
  \draw[simInk] ($(Ptip)+(-0.18,0.0)$) -- ($(Ptip)+(-0.18,0.18)$) -- ($(Ptip)+(0.0,0.18)$);
  % normalize annotation
  \node[font=\scriptsize\itshape, text=simGray, align=left, anchor=west] at (4.0,2.3)
    {\textbf{1.} subtract the shadow\\(projection onto the span)};
  \node[font=\scriptsize\itshape, text=simGray, align=left, anchor=west] at (4.0,1.3)
    {\textbf{2.} normalize $\vw'$:\\$\vq_3 = \vw'/\norm{\vw'}$};
  \node[opbox, anchor=west] at (4.0,0.3) {$\vq_3 \perp \operatorname{span}\{\vq_1,\vq_2\}$};
\end{tikzpicture}
\caption[Gram--Schmidt, geometrically]{One Gram--Schmidt step, seen in three
dimensions. The new vector $\vw$ (rust) sticks up out of the plane already spanned
by $\vq_1,\vq_2$ (green). Its \emph{shadow} on that plane is the projection
$\sum_j \inner{\vw}{\vq_j}\vq_j$ (blue)---the part of $\vw$ we already have.
Subtracting the shadow leaves the orthogonal residual $\vw'$ (violet), which by
construction meets the plane at a right angle and so points in a genuinely new
direction. Normalizing $\vw'$ to unit length yields the next basis vector $\vq_3$.
Every orthogonalization method in this chapter is some scheduling of exactly this
\emph{subtract-the-projections-then-normalize} step; what differs is the
\emph{order} in which the projections are subtracted.}
\label{fig:gschmidt}
\end{figure}

That the residual really is orthogonal to everything we kept is a one-line
calculation, and worth doing once because it is the contract the whole chapter
guarantees.

\begin{lemma}[The residual is orthogonal to the span]
\label{lem:residorthog}
Let $\vq_1,\dots,\vq_m$ be orthonormal and
$\vw' = \vw - \sum_{j=1}^m \inner{\vw}{\vq_j}\vq_j$. Then
$\inner{\vw'}{\vq_i} = 0$ for every $i \in \{1,\dots,m\}$.
\end{lemma}
\begin{proof}
Take the inner product of $\vw'$ with $\vq_i$ and use linearity:
\[
  \inner{\vw'}{\vq_i}
    = \inner{\vw}{\vq_i} - \sum_{j=1}^m \inner{\vw}{\vq_j}\,\inner{\vq_j}{\vq_i}
    = \inner{\vw}{\vq_i} - \sum_{j=1}^m \inner{\vw}{\vq_j}\,\delta_{ji}
    = \inner{\vw}{\vq_i} - \inner{\vw}{\vq_i} = 0,
\]
where the middle step used orthonormality, $\inner{\vq_j}{\vq_i} = \delta_{ji}$, to
collapse the sum to its single $j = i$ term. \qedhere
\end{proof}

The coefficients $\inner{\vw}{\vq_j}$ are not waste; they are the record of
\emph{how much} of $\vw$ lay in each previous direction. Stacked into a column they
become the entries of the \emph{Hessenberg matrix} that Arnoldi produces
(\cref{sec:arnoldi}), and the solver downstream reads them as the projection of the
operator onto the Krylov basis. Orthogonalization thus produces two things at once:
the new basis vector, and the column of coefficients. Keep both in view---they are
the $(\vw', H)$ pair that the framework's \code{orthogonalize} operator returns.

\begin{notationbox}
We follow the framework's \code{orthogonalize} contract: the operator
\emph{orthogonalizes against a normalized basis but does not itself normalize its
output.} It returns the residual $\vw'$ and the coefficient column $H$ where
$H_j = \inner{\vw}{\vq_j}$; the caller takes $\norm{\vw'}$ (which becomes the next
Hessenberg sub-diagonal entry) and divides. We keep the normalization explicit and
separate throughout, because the two recurrences of \cref{sec:arnoldi,sec:lanczos}
attach different meaning to that final norm.
\end{notationbox}

\section{Classical versus modified Gram--Schmidt}
\label{sec:cgsmgs}

The recipe of \cref{sec:gschmidt} hides a scheduling decision that does not matter
in exact arithmetic and matters enormously in floating point. The residual is a sum
of $m$ projection-subtractions; the question is \emph{against which vector each
projection coefficient is computed}. There are two natural answers, and they are the
two classical algorithms.

\subsection{Classical Gram--Schmidt (CGS)}

The literal reading of the formula
$\vw' = \vw - \sum_j \inner{\vw}{\vq_j}\vq_j$ computes every coefficient against the
\emph{same original} $\vw$:
\[
  H_j = \inner{\vw}{\vq_j} \quad\text{for all } j, \qquad
  \vw' = \vw - \sum_{j=1}^m H_j\,\vq_j .
\]
All $m$ inner products read the input $\vw$ and nothing else, so they are mutually
independent: they can be computed in any order, batched into a single matrix-vector
product $\mat{Q}^{\!\top}\vw$, and---on a parallel machine---reduced in one
collective communication. This is \emph{classical Gram--Schmidt}, CGS, and its
defining structural feature is that \emph{the $m$ projections do not depend on one
another}.

\subsection{Modified Gram--Schmidt (MGS)}

The alternative subtracts the projections \emph{one at a time}, and each coefficient
is computed against the \emph{partially-orthogonalized} vector left by the previous
subtractions. Write $\vw^{(0)} = \vw$ and, for $j = 1, \dots, m$,
\[
  H_j = \inner{\vw^{(j-1)}}{\vq_j}, \qquad
  \vw^{(j)} = \vw^{(j-1)} - H_j\,\vq_j,
\]
and set $\vw' = \vw^{(m)}$. This is \emph{modified Gram--Schmidt}, MGS. Each step
projects the \emph{running residual} onto the next basis vector, rather than
projecting the original. Equivalently, MGS applies a left-to-right composition of
rank-one projectors,
\[
  \vw' = \bigl(\mat{I} - \vq_m\vq_m^{\!\top}\bigr)\cdots
         \bigl(\mat{I} - \vq_1\vq_1^{\!\top}\bigr)\,\vw,
\]
where CGS applies the single combined projector
$\bigl(\mat{I} - \sum_j \vq_j\vq_j^{\!\top}\bigr)\vw$ in one shot.

\begin{keyidea}
CGS and MGS compute the \emph{same orthogonal residual in exact arithmetic}---both
are the projection $(\mat{I} - \mat{Q}\mat{Q}^{\!\top})\vw$, and the projection does
not care how you bracket the subtractions. They differ only in
\emph{against which vector each coefficient is taken}: CGS uses the original $\vw$
for all of them (independent, parallelizable), MGS uses the running residual
$\vw^{(j-1)}$ for the $j$-th (each depends on the last). That single difference of
\emph{ordering} is the whole story---it has no effect in exact arithmetic and a
decisive effect in finite precision, and it is the reason both algorithms exist.
\end{keyidea}

\subsection{The same residual, two schedules}

The algebraic equivalence is worth seeing directly. In MGS,
\[
  \vw^{(j)} = \vw^{(j-1)} - \inner{\vw^{(j-1)}}{\vq_j}\vq_j .
\]
But $\vw^{(j-1)} = \vw - \sum_{i<j}\inner{\vw^{(i-1)}}{\vq_i}\vq_i$ already has the
$\vq_i$ ($i<j$) components removed, and since the basis is orthonormal,
$\inner{\vw^{(j-1)}}{\vq_j} = \inner{\vw}{\vq_j}$ exactly---subtracting components
along \emph{other} basis vectors does not change the component along $\vq_j$. So in
exact arithmetic MGS's coefficient equals CGS's, the running subtraction reaches the
same place, and $\vw^{(m)} = \vw - \sum_j \inner{\vw}{\vq_j}\vq_j$. The two are one
computation with two bracketings. \Cref{fig:cgsmgs} draws the data-flow difference
that this algebraic identity papers over.

\begin{figure}[t]
\centering
\begin{tikzpicture}[font=\footnotesize]
  % ===== LEFT: CGS = fan-out, independent =====
  \begin{scope}
    \node[font=\small\bfseries, text=simBlue] at (1.4,3.0) {CGS: independent projections};
    \node[data, minimum width=10mm] (w) at (1.4,2.2) {$\vw$};
    \foreach \i/\x in {1/0,2/1.4,3/2.8}{
      \node[opbox, minimum width=11mm] (p\i) at (\x,1.0)
        {$\inner{\vw}{\vq_{\i}}$};
      \draw[flow] (w.south) -- (p\i.north);
    }
    \node[product, minimum width=30mm] (sum) at (1.4,-0.3)
      {$\vw' = \vw - \sum_j H_j\vq_j$};
    \foreach \i in {1,2,3}{\draw[flow] (p\i.south) -- (sum.north);}
    \node[font=\scriptsize, text=simGreen, align=center] at (1.4,-1.3)
      {all read the \emph{same} $\vw$:\\one batched reduction $\mat{Q}^{\!\top}\vw$,\\
       \emph{parallelizable}};
  \end{scope}
  \draw[simGray!50, thick] (4.2,-1.7) -- (4.2,3.2);
  % ===== RIGHT: MGS = sequential chain =====
  \begin{scope}[xshift=5.2cm]
    \node[font=\small\bfseries, text=simRust] at (2.3,3.0) {MGS: sequential subtraction};
    \node[data, minimum width=9mm] (w0) at (-0.4,1.0) {$\vw^{(0)}$};
    \foreach \i/\x in {1/1.0,2/2.5,3/4.0}{
      \node[opbox, minimum width=10mm] (s\i) at (\x,1.0) {$-H_{\i}\vq_{\i}$};
    }
    \node[product, minimum width=10mm] (w3) at (5.4,1.0) {$\vw'$};
    \draw[flow] (w0.east) -- (s1.west);
    \draw[flow] (s1.east) -- node[above,font=\scriptsize]{$\vw^{(1)}$} (s2.west);
    \draw[flow] (s2.east) -- node[above,font=\scriptsize]{$\vw^{(2)}$} (s3.west);
    \draw[flow] (s3.east) -- (w3.west);
    % each coefficient READS the incoming running residual
    \foreach \i/\x in {1/1.0,2/2.5,3/4.0}{
      \node[font=\scriptsize, text=simGold!80!black] at (\x,2.05)
        {$H_{\i}=\inner{\vw^{(\the\numexpr\i-1\relax)}}{\vq_{\i}}$};
      \draw[refedge, simGold!80!black] (\x,1.85) -- (s\i.north);
    }
    \node[font=\scriptsize, text=simRust, align=center] at (2.5,-0.6)
      {each $H_j$ reads the previous residual $\vw^{(j-1)}$:\\
       \textbf{column order is a sequential dependency},\\ \emph{not parallelizable}};
  \end{scope}
\end{tikzpicture}
\caption[Classical versus modified Gram--Schmidt]{The data-flow contrast that the
algebra hides. \emph{Left (CGS):} every projection coefficient $H_j$ reads the
\emph{same} original $\vw$, so the $m$ inner products fan out independently---they
are a single batched reduction $\mat{Q}^{\!\top}\vw$, reorderable and parallelizable,
with no arrows \emph{between} them. \emph{Right (MGS):} the coefficient $H_j$ reads
the \emph{running residual} $\vw^{(j-1)}$ left by the previous subtraction (gold
read-arrows), so the steps form a strict chain $\vw^{(0)}\!\to\!\vw^{(1)}\!\to\!
\cdots\!\to\!\vw'$ in which step $j$ cannot start until step $j-1$ finishes. This is
exactly the \code{map}-versus-\code{fold} split of \cref{ch:combinators}: CGS is a
\code{map} of independent dots, MGS is a \code{fold} threading the residual. The
column order is a genuine sequential dependency, the obstruction of
\cref{ch:rotations} in miniature.}
\label{fig:cgsmgs}
\end{figure}

\subsection{MGS is a sequential obstruction}

The right-hand side of \cref{fig:cgsmgs} should look familiar: it is precisely the
\code{fold} of \cref{ch:combinators}, and its carried dependency is precisely the
\emph{sequential obstruction} of \cref{ch:rotations}. MGS threads a running residual
through a chain of subtractions; the $j$-th step \emph{reads the output of the
$(j-1)$-th}. That is the loop-carried dependency that blocks the iteration rotation
($L_2 \to L_3$): you cannot turn MGS into a single ``orthogonalize against all
columns at once'' whole-field operation, because there is no such operation---the
later subtractions must see the results of the earlier ones.

\begin{pitfall}
The temptation, looking at CGS, is to ``optimize'' MGS into it: both compute the same
residual in exact arithmetic, and CGS parallelizes, so why not always use CGS? Because
the equivalence is \emph{exact-arithmetic only}. In finite precision CGS is
numerically unstable---it loses orthogonality fast for ill-conditioned bases
(\cref{sec:lossorthog})---while MGS holds it far better. The sequential dependency of
MGS is not an accident to be optimized away; it is \emph{load-bearing}. Each step's
subtraction against the freshly-cleaned residual is exactly what keeps the rounding
error from accumulating. Replacing MGS with CGS for speed silently degrades the
basis, and the degradation only shows up later, as a solver that stalls or a
spurious eigenvalue. This is the canonical case where a \emph{sequential dependency
is the algorithm}, not an obstacle to it.
\end{pitfall}

There is an algebraic fingerprint of this dependency, and it is the cleanest way to
state ``MGS does not parallelize.'' For CGS, permuting the columns of $\mat{Q}$
permutes the coefficients but leaves the residual unchanged (up to reduction-order
rounding), because the combined projector $\mat{I} - \sum_j\vq_j\vq_j^{\!\top}$ is
symmetric in its terms. For MGS, permuting the columns changes the intermediate
residuals $\vw^{(j)}$ and hence the bit-level result: the rank-one projectors
$(\mat{I} - \vq_j\vq_j^{\!\top})$ do not commute in finite precision. \emph{Column
order matters for MGS and not for CGS}---that non-commutativity is the algebraic
shadow of the sequential obstruction.

\subsection{A worked comparison on nearly parallel vectors}

The whole point of preferring MGS is invisible until we watch both run in finite
precision on an ill-conditioned input. \Cref{wex:cgsmgs} does exactly that.

\begin{workedexample}{CGS versus MGS on three nearly parallel vectors}{cgsmgs}
We orthonormalize three nearly parallel vectors in $\Real^3$ and measure how
orthogonal the resulting basis actually is. Take a small parameter
$\varepsilon = 10^{-4}$ and the columns
\[
  \va_1 = \begin{psmallmatrix} 1 \\ \varepsilon \\ 0 \end{psmallmatrix},\quad
  \va_2 = \begin{psmallmatrix} 1 \\ 0 \\ \varepsilon \end{psmallmatrix},\quad
  \va_3 = \begin{psmallmatrix} 1 \\ \varepsilon \\ \varepsilon \end{psmallmatrix}.
\]
All three are within $\varepsilon$ of the vector $(1,0,0)^{\!\top}$, so they are
\emph{nearly parallel}: a textbook ill-conditioned input, exactly the situation the
power basis produces. We carry the arithmetic as if in a precision coarse enough
that terms of order $\varepsilon^2$ relative to $1$ are lost---the qualitative
behavior is the same in IEEE double, just at $\varepsilon \approx 10^{-8}$.

\smallskip
\emph{Step 1 (shared by both).} Normalize $\va_1$. Its norm is
$\sqrt{1 + \varepsilon^2} \approx 1$, so $\vq_1 = \va_1 = (1, \varepsilon, 0)^{\!\top}$
to working precision.

\smallskip
\emph{Step 2 (shared).} Orthogonalize $\va_2$ against $\vq_1$. The coefficient is
$\inner{\va_2}{\vq_1} = 1\cdot 1 + 0\cdot\varepsilon + \varepsilon\cdot 0 = 1$, so
\[
  \vw' = \va_2 - 1\cdot\vq_1
       = \begin{psmallmatrix}1\\0\\\varepsilon\end{psmallmatrix}
       - \begin{psmallmatrix}1\\\varepsilon\\0\end{psmallmatrix}
       = \begin{psmallmatrix}0\\-\varepsilon\\\varepsilon\end{psmallmatrix},
  \qquad
  \vq_2 = \frac{\vw'}{\norm{\vw'}}
        = \frac{1}{\sqrt2}\begin{psmallmatrix}0\\-1\\1\end{psmallmatrix}.
\]
So far CGS and MGS agree---the divergence is in step 3, the first step with
\emph{two} prior vectors to project against.

\smallskip
\emph{Step 3, CGS.} Both coefficients are taken against the \emph{original} $\va_3$:
\[
  H_1 = \inner{\va_3}{\vq_1} = 1 + \varepsilon^2 \approx 1, \qquad
  H_2 = \inner{\va_3}{\vq_2}
      = \tfrac{1}{\sqrt2}(0 - \varepsilon + \varepsilon) = 0.
\]
Then
\[
  \vw'_{\text{CGS}} = \va_3 - H_1\vq_1 - H_2\vq_2
    = \begin{psmallmatrix}1\\\varepsilon\\\varepsilon\end{psmallmatrix}
    - \begin{psmallmatrix}1\\\varepsilon\\0\end{psmallmatrix}
    = \begin{psmallmatrix}0\\0\\\varepsilon\end{psmallmatrix},
  \qquad
  \vq_3^{\text{CGS}} = \begin{psmallmatrix}0\\0\\1\end{psmallmatrix}.
\]
This looks clean---but the trap is in $H_1$. The exact $\inner{\va_3}{\vq_1}$ is
$1 + \varepsilon^2$, and in finite precision the $\varepsilon^2$ term is
\emph{below the rounding threshold relative to $1$} and is dropped. The CGS residual
is computed by subtracting a rounded $H_1 \vq_1$ from $\va_3$, and the small error in
$H_1$ contaminates the $\vq_1$-direction of the residual: $\vq_3^{\text{CGS}}$ retains
a component of size $\sim\varepsilon$ along $\vq_1$ that should have been removed.
Measuring the residual orthogonality after the fact:
\[
  \inner{\vq_3^{\text{CGS}}}{\vq_1} \;\sim\; \varepsilon, \qquad
  \inner{\vq_3^{\text{CGS}}}{\vq_2} \;\sim\; \varepsilon .
\]

\smallskip
\emph{Step 3, MGS.} Now subtract sequentially. First against $\vq_1$:
\[
  H_1 = \inner{\va_3}{\vq_1} \approx 1, \qquad
  \vw^{(1)} = \va_3 - H_1\vq_1
    = \begin{psmallmatrix}0\\0\\\varepsilon\end{psmallmatrix}
    \quad(\text{to working precision}).
\]
Then take the \emph{second} coefficient against the cleaned residual $\vw^{(1)}$, not
against $\va_3$:
\[
  H_2 = \inner{\vw^{(1)}}{\vq_2}
      = \tfrac{1}{\sqrt2}(0 - 0 + \varepsilon) = \tfrac{\varepsilon}{\sqrt2},
  \qquad
  \vw^{(2)} = \vw^{(1)} - H_2\vq_2
    = \begin{psmallmatrix}0\\0\\\varepsilon\end{psmallmatrix}
      - \tfrac{\varepsilon}{2}\begin{psmallmatrix}0\\-1\\1\end{psmallmatrix}
    = \tfrac{\varepsilon}{2}\begin{psmallmatrix}0\\1\\1\end{psmallmatrix}.
\]
The crucial difference: MGS \emph{noticed} the residual's small $\vq_2$-overlap
($H_2 = \varepsilon/\sqrt2$, which CGS computed as $0$ because it looked at the
original $\va_3$ where the overlap was masked by the large shared component) and
removed it. The resulting $\vq_3^{\text{MGS}} = \tfrac{1}{\sqrt2}(0,1,1)^{\!\top}$
satisfies, to working precision,
\[
  \inner{\vq_3^{\text{MGS}}}{\vq_1} \;\sim\; \varepsilon^2, \qquad
  \inner{\vq_3^{\text{MGS}}}{\vq_2} \;\sim\; \varepsilon^2 .
\]

\smallskip
\emph{The verdict.} Both algorithms produce a unit vector roughly orthogonal to the
first two; but the \emph{residual orthogonality} differs by a full power of
$\varepsilon$: CGS leaves overlaps of size $\varepsilon$, MGS of size
$\varepsilon^2$. At $\varepsilon = 10^{-4}$ that is $10^{-4}$ versus $10^{-8}$; at the
IEEE-double analogue it is the difference between a basis that is orthogonal to four
digits and one orthogonal to eight. The reason is exactly the sequential dependency:
MGS's second projection sees a residual from which the first overlap has \emph{already}
been cleaned, so it can detect and remove the small leftover; CGS's second projection
looks at the raw input, where the same small overlap is buried under the dominant
shared component and rounds away. Order is what bought the extra digits.
\end{workedexample}

\section{Loss of orthogonality and re-orthogonalization}
\label{sec:lossorthog}

\Cref{wex:cgsmgs} showed a single step; the effect compounds. As the basis grows and
the input vectors grow more nearly parallel---which is exactly what the Krylov power
basis does---both algorithms accumulate error, and the orthogonality of the computed
basis decays. The standard way to measure the decay is the \emph{loss of
orthogonality},
\[
  \eta_m \;=\; \norm{\mat{Q}_m^{\!\top}\mat{Q}_m - \mat{I}}_2,
\]
the distance of the Gram matrix from the identity after $m$ vectors. For a perfectly
orthonormal basis $\eta_m = 0$; in finite precision it grows, and how fast it grows
is the sharpest distinction among the methods.

\begin{figure}[t]
\centering
\begin{tikzpicture}
\begin{axis}[
  width=11.5cm, height=6.6cm,
  xlabel={basis size $m$},
  ylabel={loss of orthogonality $\eta_m=\norm{\mat{Q}_m^{\!\top}\mat{Q}_m-\mat{I}}$},
  ymode=log, ymin=1e-17, ymax=1e0,
  xmin=0, xmax=30,
  grid=both, grid style={simGray!18},
  legend pos=south east, legend cell align=left,
  legend style={font=\scriptsize, fill=white, draw=simGray!40},
  label style={font=\footnotesize}, tick label style={font=\scriptsize},
  ytick={1e-16,1e-12,1e-8,1e-4,1e0},
]
% CGS: loses orthogonality quadratically in the condition number -> climbs fast to O(1)
\addplot[simRust, very thick, mark=*, mark size=1.1pt] coordinates {
  (1,2e-16)(3,5e-15)(5,8e-13)(7,4e-11)(9,9e-10)(11,3e-8)(13,6e-7)(15,2e-5)
  (17,4e-4)(19,6e-3)(21,5e-2)(23,3e-1)(25,7e-1)(27,9e-1)(29,1.0)};
\addlegendentry{CGS (unstable: $\eta\!\sim\!\kappa^2 u$)}
% MGS: loses orthogonality linearly in kappa -> climbs slowly, stays well below 1
\addplot[simBlue, very thick, mark=square*, mark size=1.1pt] coordinates {
  (1,2e-16)(3,4e-16)(5,9e-16)(7,2e-15)(9,5e-15)(11,1e-14)(13,3e-14)(15,7e-14)
  (17,2e-13)(19,5e-13)(21,1e-12)(23,3e-12)(25,7e-12)(27,1.4e-11)(29,3e-11)};
\addlegendentry{MGS (stable: $\eta\!\sim\!\kappa u$)}
% CGS2 / MGS+reorth: flat at machine epsilon
\addplot[simGreen, very thick, mark=triangle*, mark size=1.3pt] coordinates {
  (1,2e-16)(3,2e-16)(5,3e-16)(7,3e-16)(9,3e-16)(11,4e-16)(13,4e-16)(15,4e-16)
  (17,5e-16)(19,5e-16)(21,5e-16)(23,5e-16)(25,6e-16)(27,6e-16)(29,6e-16)};
\addlegendentry{CGS2 / re-orthogonalized (flat at $u$)}
% machine epsilon reference line
\addplot[simGray, densely dashed, thick, domain=0:30, samples=2] {2.2e-16};
\node[font=\scriptsize, text=simGray, anchor=west] at (axis cs:0.5,5e-16) {machine $u$};
\end{axis}
\end{tikzpicture}
\caption[Loss of orthogonality versus iteration]{How orthogonality decays as the
basis grows, on an ill-conditioned (nearly rank-deficient) input---a schematic of the
classic experiment. The vertical axis is the loss of orthogonality
$\eta_m = \norm{\mat{Q}_m^{\!\top}\mat{Q}_m - \mat{I}}$, log-scaled. \emph{CGS}
(rust) loses orthogonality \emph{quadratically} in the condition number,
$\eta\sim\kappa^2 u$, climbing all the way to $O(1)$---a basis that is no longer
orthogonal at all. \emph{MGS} (blue) loses it only \emph{linearly},
$\eta\sim\kappa u$, staying many orders of magnitude better for the same input.
\emph{CGS2}---classical Gram--Schmidt done twice---and MGS-with-one-reorthogonalization
(green) stay pinned at machine precision $u$ regardless of $m$: full orthogonality,
recovered. The flat green line is the ``twice is enough'' result of
\cref{sec:twiceenough} made visible.}
\label{fig:lossorthog}
\end{figure}

\Cref{fig:lossorthog} is the picture to carry away. On an ill-conditioned input the
three regimes separate cleanly. CGS loses orthogonality \emph{quadratically} in the
condition number $\kappa$ of the input: $\eta \sim \kappa^2 u$, where $u$ is the unit
roundoff. For an input with $\kappa = 10^8$ in double precision
($u \approx 10^{-16}$), that is $\eta \sim 10^{16}\cdot 10^{-16} = 1$---total loss of
orthogonality, exactly what \cref{wex:cgsmgs} foreshadowed compounding over many
steps. MGS loses it only \emph{linearly}, $\eta \sim \kappa u$, which for the same
input is $\eta \sim 10^{-8}$: degraded, but usable. The extra factor of $\kappa$ that
CGS pays is the price of computing every coefficient against the contaminated
original instead of the cleaned residual.

\begin{pitfall}
Even MGS is not safe for the worst inputs. Its loss $\eta \sim \kappa u$ is small only
while $\kappa$ is moderate; for a sufficiently ill-conditioned basis---and a Krylov
power basis becomes arbitrarily ill-conditioned as $m$ grows---$\kappa u$ can itself
reach $O(1)$, and MGS too produces a basis that is no longer orthogonal. The blue
curve in \cref{fig:lossorthog} stays low because the example's $\kappa$ is moderate;
push $\kappa$ high enough and it climbs. The lesson is that \emph{neither} classical
nor modified Gram--Schmidt is unconditionally trustworthy. When orthogonality must be
guaranteed to machine precision---and for Krylov eigensolvers it must---one needs the
re-orthogonalization of the next subsection.
\end{pitfall}

\subsection{``Twice is enough'': CGS2 and re-orthogonalization}
\label{sec:twiceenough}

The repair is disarmingly simple: \emph{do it again}. After computing a residual
$\vw'$ that should be orthogonal to the basis but, in finite precision, retains a
small overlap, run the \emph{whole} Gram--Schmidt step a second time on $\vw'$. The
second pass projects out the small residual overlap that the first pass left behind.
For classical Gram--Schmidt this two-pass version is called \emph{CGS2}; the same
idea applied to MGS is \emph{MGS with re-orthogonalization}.

The remarkable fact, due to Kahan and analyzed by Parlett and others, is that
\emph{one} extra pass suffices: after re-orthogonalizing once, the basis is
orthogonal to machine precision, and a third pass buys nothing. This is the ``twice
is enough'' result, and it is what flattens the green curve in \cref{fig:lossorthog}
at $u$.

\begin{figure}[t]
\centering
\begin{tikzpicture}[font=\footnotesize, scale=1.0, >={Stealth[length=2.4mm]}]
  % the span (a horizontal line in 2D representing span{q1,...})
  \draw[very thick, simGreen] (-0.3,0) -- (5.5,0);
  \node[font=\scriptsize\itshape, text=simGreen, anchor=west] at (4.6,-0.35)
    {$\operatorname{span}\{\vq_1,\dots\}$};
  \coordinate (O) at (0.4,0);
  % input w (well off the span)
  \draw[flow,simRust,very thick] (O) -- (1.0,2.4) node[above,font=\scriptsize]{$\vw$};
  % first pass result: SHOULD be vertical but tilts slightly (residual overlap)
  \coordinate (w1) at (2.6,2.3);
  \draw[flow,simBlue,very thick] (O) -- (w1)
    node[right,font=\scriptsize]{$\vw'$ (pass 1)};
  \node[font=\scriptsize,text=simBlue] at (3.1,1.4) {small residual overlap $\swarrow$};
  % the leftover horizontal component of w1
  \draw[densely dashed, simGray] (w1) -- (2.6,0);
  \draw[<->, simRust, thick] (0.4,-0.55) -- (2.6,-0.55)
    node[midway,below,font=\scriptsize,text=simRust]{$\sim u\kappa$ overlap remains};
  % second pass result: truly vertical
  \coordinate (w2) at (0.4,2.3);
  \draw[flow,simGreen,very thick] (O) -- (w2)
    node[left,font=\scriptsize]{$\vw''$ (pass 2)};
  \draw[simInk] (0.4,0.0) ++(0,0.2) -- ++(0.2,0) -- ++(0,-0.2);
  \node[font=\scriptsize,text=simGreen,anchor=west] at (0.7,2.0)
    {orthogonal to $\sim u$};
  % the correction arrow from w1 to w2
  \draw[->, simViolet, densely dashed, thick] (w1) -- (w2)
    node[midway,above,font=\scriptsize,text=simViolet]{2nd projection removes it};
\end{tikzpicture}
\caption[``Twice is enough'' re-orthogonalization]{Why one extra pass restores
orthogonality. The input $\vw$ (rust) is orthogonalized against the span (green
line). In finite precision the first pass produces $\vw'$ (blue) that is \emph{nearly}
perpendicular but retains a small residual overlap with the span---of relative size
$\sim u\kappa$, the loss of orthogonality. Running the \emph{same} Gram--Schmidt step
a second time on $\vw'$ projects out that small leftover (violet correction),
yielding $\vw''$ (green) orthogonal to machine precision $u$. The Kahan/Parlett
result is that this happens in \emph{one} re-pass: the second residual overlap is of
size $\sim u\cdot u\kappa$, already negligible, so a third pass is unnecessary. Twice
is enough.}
\label{fig:twice}
\end{figure}

\begin{notationbox}
CGS2 returns the \emph{combined} coefficients. The first pass computes a coefficient
column $H$ and a residual $\vw'$; the second pass computes a small \emph{correction}
column $\Delta H$ against $\vw'$ and a twice-projected residual $\vw''$. The operator
returns $H + \Delta H$ (so that $\vw = \vw'' + \sum_j (H_j + \Delta H_j)\vq_j$ still
recovers the input exactly) and the residual $\vw''$. The second pass is \emph{not}
algebraically fusible into the first---it reads the once-orthogonalized $\vw'$, which
did not exist until the first pass ran---so CGS2 genuinely costs two passes of inner
products. That doubled cost, traded against $\kappa^2 u \to u$ orthogonality, is the
CGS2 bargain.
\end{notationbox}

The cost/stability trade is now a clean menu, and it is worth tabulating because the
solver chooses from it at configuration time (\cref{tab:gstradeoff}).

\begin{table}[t]
\centering
\caption[The orthogonalization cost/stability trade]{The orthogonalization menu. The
\emph{communication} column counts global reductions per orthogonalization against an
$m$-vector basis (the cost that dominates on a parallel machine): CGS batches all $m$
inner products into one reduction, MGS needs $m$ separate ones because each waits on
the previous subtraction. The \emph{orthogonality} column is the resulting loss
$\eta\sim\norm{\mat{Q}^{\!\top}\mat{Q}-\mat{I}}$ in terms of the condition number
$\kappa$ and unit roundoff $u$. There is no free lunch: the cheapest-to-communicate
method (CGS) is the least stable, and the most stable (CGS2) pays double.}
\label{tab:gstradeoff}
\small
\setlength{\tabcolsep}{6pt}
\renewcommand{\arraystretch}{1.25}
\begin{tabular}{@{}llll@{}}
\toprule
Method & Passes & Communication & Orthogonality $\eta$ \\
\midrule
CGS  & 1 & 1 reduction (batched) & $\sim \kappa^2 u$ \quad(unstable) \\
MGS  & 1 & $m$ reductions (sequential) & $\sim \kappa\, u$ \quad(stable) \\
CGS2 & 2 & 2 reductions (batched) & $\sim u$ \quad(full) \\
MGS+reorth & 2 & $2m$ reductions & $\sim u$ \quad(full) \\
\bottomrule
\end{tabular}
\end{table}

CGS2 occupies a sweet spot the table makes plain: it recovers full orthogonality
($\eta\sim u$) like MGS-with-reorthogonalization, but because each of its two passes
\emph{batches} its $m$ inner products into a single reduction, it pays only
\emph{two} communications rather than $2m$. On a parallel machine where communication
dominates, CGS2 is often the method of choice precisely because it buys MGS-quality
stability while keeping CGS-style batched communication. The sequential dependency
that made MGS stable is also what makes it communication-bound; CGS2 sidesteps both
by re-running the parallel-friendly classical pass.

\section{Arnoldi: Gram--Schmidt on the Krylov sequence}
\label{sec:arnoldi}

We now assemble the orthogonalizer into the recurrence that drives GMRES. The
\emph{Arnoldi iteration} is nothing more than Gram--Schmidt applied to the Krylov
sequence as it is generated: start from a unit vector $\vq_1$, and at each step form
the next power-basis vector $\mat{A}\vq_j$, orthogonalize it against all the basis
vectors so far, and normalize the residual to get $\vq_{j+1}$.

\begin{lstlisting}[style=l4style]
arnoldi_step :: LinOp -> Basis -> Tensor[$S] -> (Tensor[$S], Tensor[$m+1])
arnoldi_step A V q_j =
  let w        = apply A q_j                    -- next power-basis vector A q_j
      (w', h)  = orthogonalize w V CGS2         -- project off span(V); h = coeffs H[1..m]
      beta     = nrm2 w'                         -- the sub-diagonal entry h_{m+1,m}
      q_next   = scal (1 / beta) w'              -- normalize: the new basis vector
  in (q_next, append h beta)                     -- (new column of V, full Hessenberg column)
\end{lstlisting}

The coefficients produced along the way are not discarded---they are the entries of a
matrix. After $m$ steps, the basis vectors $\vq_1,\dots,\vq_m$ form the columns of an
$N\times m$ matrix $\mat{V}_m$, and the orthogonalization coefficients fill an
$(m+1)\times m$ matrix $\bar{\mat{H}}_m$ that is \emph{upper Hessenberg}: zero below
the first sub-diagonal. The entry $H_{ij}$ is the coefficient of $\vq_i$ when
$\mat{A}\vq_j$ was orthogonalized---which is zero whenever $i > j+1$, because at step
$j$ there were only $j$ basis vectors to project against plus the one new
sub-diagonal normalization. These fit together into the \emph{Arnoldi relation}
\[
  \mat{A}\,\mat{V}_m \;=\; \mat{V}_{m+1}\,\bar{\mat{H}}_m,
\]
which says: applying $\mat{A}$ to the basis and re-expressing the result in the basis
\emph{is} the Hessenberg matrix. \Cref{fig:arnoldi} draws the relation and the shapes
involved.

\begin{figure}[t]
\centering
\begin{tikzpicture}[font=\footnotesize, scale=0.9]
  % A V_m
  \node[font=\small] at (-5.4,0) {$\mat{A}\,\mat{V}_m$};
  \node[font=\small] at (-4.3,0) {$=$};
  % V_{m+1} (tall thin block)
  \draw[thick, fill=simBgGreen] (-3.7,-1.5) rectangle (-2.5,1.5);
  \foreach \x in {-3.45,-3.2,-2.95,-2.7}{\draw[simGreen!60] (\x,-1.5)--(\x,1.5);}
  \node[font=\small] at (-3.1,-2.0) {$\mat{V}_{m+1}$};
  \node[font=\scriptsize,text=simGray] at (-3.1,1.85) {$N\times(m{+}1)$};
  \node[font=\small] at (-2.15,0) {$\cdot$};
  % H-bar (upper Hessenberg: staircase)
  \draw[thick, fill=simBgBlue] (-1.6,-1.5) rectangle (0.4,1.5);
  % shade the zero (below sub-diagonal) region
  \fill[white] (-1.6,-1.5) -- (0.4,1.5) ++(0,-0.6) -- (0.4,-1.5) -- cycle;
  % redraw the staircase boundary of the nonzero (upper Hessenberg) part
  \draw[thick, simBlue]
    (-1.6,1.5) -- (0.4,1.5) -- (0.4,0.9)
    -- (0.0,0.9) -- (0.0,0.3) -- (-0.4,0.3) -- (-0.4,-0.3)
    -- (-0.8,-0.3) -- (-0.8,-0.9) -- (-1.2,-0.9) -- (-1.2,-1.5) -- (-1.6,-1.5) -- cycle;
  \node[font=\small] at (-0.6,-2.0) {$\bar{\mat{H}}_m$};
  \node[font=\scriptsize,text=simGray] at (-0.6,1.85) {$(m{+}1)\times m$};
  \node[font=\scriptsize,text=simBlue] at (-0.55,0.55) {upper};
  \node[font=\scriptsize,text=simBlue] at (-0.55,0.15) {Hessenberg};
  \node[font=\scriptsize,text=simGray] at (-0.05,-1.15) {sub-diag};
  \node[font=\scriptsize,text=simGray] at (-1.15,-0.55) {$=\norm{\vw'}$};
  % annotation on the right
  \node[align=left, font=\scriptsize, text=simGray, anchor=west] at (1.2,0.6)
    {$\bar{\mat{H}}_m[i,j]$ = coefficient of $\vq_i$\\
     when $\mat{A}\vq_j$ was orthogonalized;\\
     the \emph{projection coefficients} $H$\\
     of \cref{sec:gschmidt} stacked column by column.};
  \node[align=left, font=\scriptsize, text=simGreen, anchor=west] at (1.2,-0.7)
    {sub-diagonal $\bar{\mat{H}}_m[j{+}1,j]=\norm{\vw'}$\\
     is the caller's \code{nrm2} normalization;\\
     zeros below it $\Rightarrow$ Hessenberg shape.};
\end{tikzpicture}
\caption[Arnoldi produces $\mat{V}$ and an upper-Hessenberg $\bar{\mat{H}}$]{The
Arnoldi relation $\mat{A}\mat{V}_m = \mat{V}_{m+1}\bar{\mat{H}}_m$, the output of
running Gram--Schmidt on the Krylov sequence. The basis vectors fill the tall thin
$\mat{V}_{m+1}$ (green); the orthogonalization coefficients fill the small
\emph{upper-Hessenberg} $\bar{\mat{H}}_m$ (blue)---zero below the first sub-diagonal,
because at step $j$ only the first $j$ basis vectors and one new normalization
contribute. Each column of $\bar{\mat{H}}_m$ is exactly the $(\vw', H)$ coefficient
pair of \cref{sec:gschmidt}: the entries $H_{ij}=\inner{\mat{A}\vq_j}{\vq_i}$ above
the diagonal, and the sub-diagonal entry the residual norm $\norm{\vw'}$ from the
caller's normalization. GMRES (\cref{ch:gmres}) reads this small Hessenberg matrix to
solve its least-squares problem; the Arnoldi eigensolvers of \cref{ch:krylovschur}
read its eigenvalues. The orthogonalizer of this chapter is the engine that fills
both matrices.}
\label{fig:arnoldi}
\end{figure}

This is precisely the \code{orthogonalize} step GMRES (\cref{ch:gmres}) calls once
per inner iteration, and the upper-Hessenberg matrix it fills is the small dense
matrix on which GMRES solves its least-squares problem. The Arnoldi eigensolvers of
\cref{ch:krylovschur} use the very same recurrence and read $\bar{\mat{H}}_m$'s
\emph{eigenvalues} (the Ritz values) as approximations to $\mat{A}$'s. Either way,
the orthogonalizer is the engine: feed it the Krylov sequence, and it produces both
the orthonormal basis and the Hessenberg matrix that is the projection of $\mat{A}$
onto that basis.

\section{Lanczos: the three-term recurrence for symmetric operators}
\label{sec:lanczos}

Now specialize to a \emph{symmetric} operator, $\mat{A} = \mat{A}^{\!\top}$ (the case
of the eigenmode analysis and of conjugate gradients). Something remarkable happens
to Arnoldi: almost all of the orthogonalization coefficients turn out to be zero in
advance, and the long Gram--Schmidt sum collapses to a \emph{three-term recurrence}.

\subsection{Why symmetry shortens the recurrence}

The Hessenberg matrix $\bar{\mat{H}}_m$ of Arnoldi is the projection of $\mat{A}$ onto
the Krylov basis: $H_{ij} = \inner{\mat{A}\vq_j}{\vq_i}$. If $\mat{A}$ is symmetric,
then so is this projection,
\[
  H_{ij} = \inner{\mat{A}\vq_j}{\vq_i} = \inner{\vq_j}{\mat{A}\vq_i}
         = \inner{\mat{A}\vq_i}{\vq_j} = H_{ji},
\]
where the middle equality is the symmetry of $\mat{A}$. But $\bar{\mat{H}}_m$ is
\emph{already} upper Hessenberg---zero below the first sub-diagonal. A matrix that is
both upper Hessenberg \emph{and} symmetric is \emph{tridiagonal}: zero everywhere
except the diagonal and the two adjacent off-diagonals. So in the symmetric case the
Hessenberg matrix is a tridiagonal matrix $\mat{T}_m$, and the Arnoldi
orthogonalization of $\mat{A}\vq_j$ has \emph{only three nonzero coefficients}: against
$\vq_{j-1}$, against $\vq_j$, and the new sub-diagonal normalization for $\vq_{j+1}$.
Every projection against $\vq_1,\dots,\vq_{j-2}$ is zero, automatically.

\begin{keyidea}
For a symmetric operator $\mat{A}$, the Arnoldi Hessenberg matrix is symmetric, and a
symmetric Hessenberg matrix is \emph{tridiagonal}. So orthogonalizing
$\mat{A}\vq_j$---which in general requires projecting against \emph{all} $j$ previous
basis vectors---requires projecting against only the \emph{two} most recent ones,
$\vq_j$ and $\vq_{j-1}$; the projections against all earlier vectors are exactly zero.
The long Gram--Schmidt sum collapses to a three-term recurrence. This is the same
miracle that gives conjugate gradients its short recurrence (\cref{ch:cg}): symmetry
makes a quantity that would otherwise depend on the whole history depend on only the
last two steps.
\end{keyidea}

The recurrence is the \emph{Lanczos iteration}. Writing the three coefficients as
$\alpha_j$ (diagonal), $\beta_{j-1}$ (the sub-diagonal coupling to the previous
vector), and $\beta_j$ (the new normalization), the step is
\[
  \vw \;=\; \mat{A}\vq_j - \alpha_j\,\vq_j - \beta_{j-1}\,\vq_{j-1},
  \qquad
  \alpha_j = \inner{\mat{A}\vq_j}{\vq_j},
  \qquad
  \beta_j = \norm{\vw},
  \qquad
  \vq_{j+1} = \frac{\vw}{\beta_j}.
\]
Only $\vq_j$ and $\vq_{j-1}$ appear---the two previous vectors, nothing earlier. In
code:

\begin{lstlisting}[style=l4style]
lanczos_step :: LinOp -> Tensor[$S] -> Tensor[$S] -> Scalar
                      -> (Tensor[$S], Scalar, Scalar)
lanczos_step A q_j q_prev beta_prev =
  let u      = apply A q_j                       -- A q_j
      alpha  = dot u q_j                          -- diagonal coeff <A q_j, q_j>
      w      = axpby (-alpha) q_j 1 u             -- subtract the q_j component
      w'     = axpy (-beta_prev) q_prev w         -- subtract the q_{j-1} component (only TWO!)
      beta   = nrm2 w'                            -- new sub-diagonal / normalization
      q_next = scal (1/beta) w'
  in (q_next, alpha, beta)
\end{lstlisting}

The contrast with \code{arnoldi\_step} is the whole point: Arnoldi's
\code{orthogonalize w V} projects against the \emph{entire} basis \code{V}, a sum
that grows with $m$; Lanczos projects against exactly two stored vectors,
$\vq_j$ and $\vq_{j-1}$, a fixed-size step. \Cref{fig:lanczos} draws the difference,
and the tridiagonal matrix the recurrence fills.

\begin{figure}[t]
\centering
\begin{tikzpicture}[font=\footnotesize, scale=0.95]
  % ===== LEFT: the three-term recurrence, only two previous vectors =====
  \begin{scope}
    \node[font=\small\bfseries, text=simGreen] at (1.4,2.9) {three-term recurrence};
    % the stored vectors
    \node[data, minimum width=11mm] (qm1) at (-0.4,1.6) {$\vq_{j-1}$};
    \node[data, minimum width=11mm] (qj)  at (1.4,1.6) {$\vq_j$};
    \node[stage, minimum width=14mm] (Aqj) at (1.4,0.5) {$\mat{A}\vq_j$};
    \draw[flow] (qj) -- (Aqj);
    \node[product, minimum width=11mm] (qn) at (1.4,-0.8) {$\vq_{j+1}$};
    % only the two recent vectors feed the new one
    \draw[flow,simGreen] (qj.south west) .. controls +(-0.2,-0.6) .. (qn.north)
      node[pos=0.5,left,font=\scriptsize,text=simGreen]{$-\alpha_j$};
    \draw[flow,simGreen] (qm1.south) .. controls +(0,-1.4) .. (qn.west)
      node[pos=0.3,left,font=\scriptsize,text=simGreen]{$-\beta_{j-1}$};
    \draw[flow] (Aqj) -- (qn);
    % the OLD vectors that are NOT needed
    \node[data, minimum width=9mm, draw=simGray!40, text=simGray!60] (old) at (-1.7,1.6) {$\vq_{j-2}\dots$};
    \draw[refedge, simRust] (old.south) .. controls +(0,-1.8) .. (qn.north west);
    \node[font=\scriptsize, text=simRust] at (-1.5,-0.5) {coeff $=0$:};
    \node[font=\scriptsize, text=simRust] at (-1.5,-0.85) {not needed};
    \node[font=\scriptsize\itshape, text=simGray, align=center] at (0.4,-1.7)
      {only $\vq_j,\vq_{j-1}$ stored:\\$O(1)$ memory per step};
  \end{scope}
  \draw[simGray!50, thick] (3.2,-1.9) -- (3.2,2.9);
  % ===== RIGHT: the tridiagonal T =====
  \begin{scope}[xshift=5.0cm]
    \node[font=\small\bfseries, text=simBlue] at (1.5,2.9) {tridiagonal $\mat{T}_m$};
    % draw a 5x5 tridiagonal pattern
    \def\cell{0.55}
    \foreach \r in {0,1,2,3,4}{
      \foreach \c in {0,1,2,3,4}{
        \pgfmathtruncatemacro\diff{abs(\r-\c)}
        \ifnum\diff<2
          \ifnum\r=\c
            \node[fill=simBgGold, draw=simGold, minimum size=\cell cm, font=\scriptsize]
              at (\c*\cell, -\r*\cell) {$\alpha_{\the\numexpr\r+1\relax}$};
          \else
            \node[fill=simBgBlue, draw=simBlue, minimum size=\cell cm, font=\scriptsize]
              at (\c*\cell, -\r*\cell) {$\beta$};
          \fi
        \else
          \node[draw=simGray!25, minimum size=\cell cm, font=\scriptsize, text=simGray!50]
            at (\c*\cell, -\r*\cell) {$0$};
        \fi
      }
    }
    \node[font=\scriptsize\itshape, text=simGray, align=center] at (1.1,-3.1)
      {symmetric + Hessenberg\\$=$ tridiagonal: $\alpha$ diag, $\beta$ off-diag};
  \end{scope}
\end{tikzpicture}
\caption[The Lanczos three-term recurrence and its tridiagonal matrix]{The Lanczos
specialization of Arnoldi for symmetric $\mat{A}$. \emph{Left:} the new basis vector
$\vq_{j+1}$ is built from $\mat{A}\vq_j$ by subtracting components along only the
\emph{two} most recent vectors, $\vq_j$ (coefficient $-\alpha_j$) and $\vq_{j-1}$
(coefficient $-\beta_{j-1}$). All earlier vectors $\vq_{j-2},\dots$ have coefficient
exactly zero (rust, ``not needed''), so they need not be stored: Lanczos uses $O(1)$
memory per step where Arnoldi stores the whole growing basis. \emph{Right:} the
coefficients fill a \emph{tridiagonal} matrix $\mat{T}_m$---the diagonal carries the
$\alpha_j = \inner{\mat{A}\vq_j}{\vq_j}$, the two off-diagonals carry the
$\beta_j = \norm{\vw}$, and everything else is zero. This is Arnoldi's Hessenberg
matrix made tridiagonal by symmetry. The short recurrence is the same economy
symmetry buys conjugate gradients (\cref{ch:cg}).}
\label{fig:lanczos}
\end{figure}

\subsection{A worked Lanczos: two steps on a small symmetric matrix}

To make the recurrence concrete, run it by hand on a $3\times 3$ symmetric matrix and
watch the tridiagonal $\mat{T}$ assemble.

\begin{workedexample}{Two Lanczos steps building a tridiagonal $\mat{T}$}{lanczos}
Take the symmetric matrix and starting vector
\[
  \mat{A} = \begin{psmallmatrix} 2 & 1 & 0 \\ 1 & 2 & 1 \\ 0 & 1 & 2 \end{psmallmatrix},
  \qquad
  \vq_1 = \begin{psmallmatrix} 1 \\ 0 \\ 0 \end{psmallmatrix}
  \quad(\text{already unit length}).
\]

\emph{Step 1.} Form $\mat{A}\vq_1$, the diagonal coefficient $\alpha_1$, and the
residual:
\[
  \mat{A}\vq_1 = \begin{psmallmatrix} 2 \\ 1 \\ 0 \end{psmallmatrix},
  \qquad
  \alpha_1 = \inner{\mat{A}\vq_1}{\vq_1} = 2 .
\]
There is no $\vq_0$, so the recurrence has only the $\alpha_1$ term on the first step:
\[
  \vw = \mat{A}\vq_1 - \alpha_1\vq_1
      = \begin{psmallmatrix} 2 \\ 1 \\ 0 \end{psmallmatrix}
      - 2\begin{psmallmatrix} 1 \\ 0 \\ 0 \end{psmallmatrix}
      = \begin{psmallmatrix} 0 \\ 1 \\ 0 \end{psmallmatrix},
  \qquad
  \beta_1 = \norm{\vw} = 1,
  \qquad
  \vq_2 = \frac{\vw}{\beta_1} = \begin{psmallmatrix} 0 \\ 1 \\ 0 \end{psmallmatrix}.
\]
Check orthogonality: $\inner{\vq_2}{\vq_1} = 0$. Good.

\emph{Step 2.} Now the full three-term step, using $\vq_2$ and the previous $\vq_1$
with its coupling $\beta_1 = 1$:
\[
  \mat{A}\vq_2 = \begin{psmallmatrix} 1 \\ 2 \\ 1 \end{psmallmatrix},
  \qquad
  \alpha_2 = \inner{\mat{A}\vq_2}{\vq_2} = 2 .
\]
Subtract \emph{both} the $\vq_2$ component (coefficient $\alpha_2$) and the $\vq_1$
component (coefficient $\beta_1$):
\[
  \vw = \mat{A}\vq_2 - \alpha_2\vq_2 - \beta_1\vq_1
      = \begin{psmallmatrix} 1 \\ 2 \\ 1 \end{psmallmatrix}
      - 2\begin{psmallmatrix} 0 \\ 1 \\ 0 \end{psmallmatrix}
      - 1\begin{psmallmatrix} 1 \\ 0 \\ 0 \end{psmallmatrix}
      = \begin{psmallmatrix} 0 \\ 0 \\ 1 \end{psmallmatrix},
  \qquad
  \beta_2 = \norm{\vw} = 1,
  \qquad
  \vq_3 = \begin{psmallmatrix} 0 \\ 0 \\ 1 \end{psmallmatrix}.
\]
Check: $\inner{\vq_3}{\vq_2} = 0$ and $\inner{\vq_3}{\vq_1} = 0$---the second of these
\emph{automatically}, even though we never projected against $\vq_1$ in step 2. That
is the symmetry miracle in action: the $\vq_1$ component of $\mat{A}\vq_2$ was already
exactly accounted for by the $\beta_1\vq_1$ term, so no further projection against
$\vq_1$ (or any earlier vector) was needed.

\emph{The assembled matrix.} The coefficients fill the tridiagonal
\[
  \mat{T}_2 = \begin{psmallmatrix} \alpha_1 & \beta_1 \\ \beta_1 & \alpha_2 \end{psmallmatrix}
            = \begin{psmallmatrix} 2 & 1 \\ 1 & 2 \end{psmallmatrix},
\]
and after the third step $\mat{T}_3 = \mathrm{tridiag}(\beta;\,\alpha;\,\beta)$ with
$\alpha = (2,2,2)$, $\beta = (1,1)$---which for this particular $\mat{A}$ is just
$\mat{A}$ itself expressed in the Lanczos basis, since the Krylov subspace happened to
fill all of $\Real^3$ in three steps. The eigenvalues of $\mat{T}_3$ are therefore
exactly those of $\mat{A}$, namely $2$ and $2\pm\sqrt2$---which is the basis on which
the eigensolvers of \cref{ch:krylovschur} read approximate eigenvalues off the small
tridiagonal $\mat{T}_m$ for $m < N$, where the agreement is approximate rather than
exact.
\end{workedexample}

\subsection{The catastrophe: Lanczos loses orthogonality}

The three-term recurrence is beautiful and economical, and in finite precision it is
treacherous. The short recurrence is short precisely because it \emph{trusts} that
the projections against $\vq_1,\dots,\vq_{j-2}$ are zero and so skips them. In exact
arithmetic they are zero. In floating point they are not---and because Lanczos never
re-checks them, the small errors are never corrected. Orthogonality against the
\emph{old} vectors is lost, and lost fast.

\begin{pitfall}
Lanczos loses orthogonality far more violently than Gram--Schmidt, and for a reason
peculiar to the short recurrence: it never looks at the old basis vectors, so it
cannot notice (let alone repair) the orthogonality it loses against them. The
practical symptom is \emph{ghost eigenvalues}. As orthogonality decays, the Lanczos
basis silently develops near-duplicate directions, and the tridiagonal $\mat{T}_m$
reports the \emph{same} eigenvalue several times over---spurious copies of converged
eigenvalues that are artifacts of the lost orthogonality, not real eigenvalues of
$\mat{A}$. A naive Lanczos eigensolver will confidently return a multiplicity-three
eigenvalue that has multiplicity one. The cure is a re-orthogonalization
\emph{policy}---full re-orthogonalization against the whole stored basis (expensive,
defeating the $O(1)$-memory point), or selective/partial re-orthogonalization that
re-orthogonalizes only when a monitored loss-of-orthogonality estimate crosses a
threshold. Which policy a solver chooses, and how the practical Krylov--Schur method
of \cref{ch:krylovschur} manages it, is a central design question of the eigensolver.
\end{pitfall}

This is the trade Lanczos forces. Arnoldi, with its full orthogonalization, is
expensive (storing and projecting against the whole basis) but stable. Lanczos, with
its three-term recurrence, is cheap ($O(1)$ memory, two projections) but unstable, and
buying back the stability with re-orthogonalization gives back some of the cost it
saved. The eigensolver chapters live inside exactly this tension.

\section{The orthogonality contract}
\label{sec:contract}

Step back from the four algorithms and notice what they all \emph{deliver}. CGS, MGS,
CGS2, Arnoldi, Lanczos---each is a different schedule, a different cost, a different
stability, a different memory footprint. But every one of them is in the business of
producing the same thing: a basis that satisfies
\[
  \inner{\vq_i}{\vq_j} \approx \delta_{ij}
  \qquad\text{to a stated tolerance,}
\]
and a coefficient matrix ($\bar{\mat{H}}$ or $\mat{T}$) that records the projection of
the operator onto that basis. This is the \emph{orthogonality contract}, and it is the
single interface every downstream method programs against.

\begin{keyidea}
\textbf{The orthogonality contract.} Every orthogonalization method in this chapter
guarantees one thing: an orthonormal basis $\mat{Q}$ with
$\mat{Q}^{\!\top}\mat{Q} \approx \mat{I}$ to a tolerance, together with the
coefficients of the projection it performed. The methods differ \emph{only} in how
well they hold the contract and at what cost---CGS cheaply but loosely, MGS stably but
sequentially, CGS2 fully but at double cost, Lanczos cheaply but fragilely. They are
therefore \emph{substitutable}: a caller that depends only on the contract (``give me
an orthonormal basis'') can swap one for another without changing its own code. This
is exactly the algorithmic-substitution that \cref{ch:rotations} calls a genuine
rotation---one slot, \lstinline[style=l4style]@orthog := mgs | cgs | cgs2@, three
different algorithms behind one interface.
\end{keyidea}

That substitutability is not a coincidence of notation; it is the contract doing its
job. GMRES does not care \emph{how} its basis was orthogonalized---it cares that the
basis is orthonormal so that its least-squares problem is well-conditioned. An Arnoldi
eigensolver does not care which Gram--Schmidt variant filled the Hessenberg
matrix---it cares that the Ritz values it reads are not corrupted by a basis that
secretly lost orthogonality. The contract is what lets the solver authors reason about
their algorithm \emph{above} the orthogonalizer, treating it as a black box with a
guarantee, while the orthogonalizer authors reason about stability and cost
\emph{below} the interface. And it is what makes ``did orthogonality hold?'' the first
question to ask when a Krylov method misbehaves: a stalled GMRES or a ghost eigenvalue
is, far more often than not, a contract quietly violated by a basis that lost its
orthogonality and a method that never noticed.

\summarysection
A Krylov solver works inside $\mathcal{K}_m(\mat{A},\vv)$ but never in its natural
power basis, which collapses toward the dominant eigenvector and is hopelessly
ill-conditioned; it works in an \emph{orthonormal} basis, built incrementally by
\emph{Gram--Schmidt}: subtract off the projections onto the previous basis vectors,
then normalize. Gram--Schmidt has two faces. \emph{Classical} (CGS) computes every
projection coefficient against the same original vector---independent, batchable,
parallelizable, but numerically unstable, losing orthogonality as $\kappa^2 u$.
\emph{Modified} (MGS) subtracts the projections one at a time against the running
residual---stable, losing orthogonality only as $\kappa u$, but at the price of a
genuine \emph{sequential dependency}: the column order is a loop-carried dependency, a
concrete specimen of the obstruction of \cref{ch:rotations}, and exactly the
\code{map}-versus-\code{fold} split of \cref{ch:combinators}. For the worst inputs
even MGS is not enough, and \emph{re-orthogonalization} (CGS2, ``twice is enough'') one
extra pass restores full machine-precision orthogonality. Applied to the Krylov
sequence, Gram--Schmidt becomes \emph{Arnoldi}, producing the orthonormal $\mat{V}_m$
and the upper-Hessenberg $\bar{\mat{H}}_m$ that GMRES (\cref{ch:gmres}) and the
Arnoldi eigensolvers (\cref{ch:krylovschur}) consume. For \emph{symmetric} operators
Arnoldi collapses to the \emph{Lanczos} three-term recurrence---the Hessenberg matrix
becomes tridiagonal, the orthogonalization needs only the two most recent vectors, the
same short-recurrence miracle that powers conjugate gradients (\cref{ch:cg})---but in
finite precision Lanczos loses orthogonality catastrophically, producing ghost
eigenvalues and demanding a re-orthogonalization policy. Underneath all of it is one
\emph{orthogonality contract}: an orthonormal basis to a tolerance, plus the
projection coefficients. The methods differ only in how well and how cheaply they hold
it, which makes them substitutable behind one interface---the genuine rotation of
\cref{ch:rotations}.

\furtherreading
The definitive textbook development of orthogonalization, conditioning, and the
stability of Gram--Schmidt is Trefethen and Bau's \emph{Numerical Linear
Algebra}~\citep{trefethenbau1997}, whose lectures on QR factorization and on the
sensitivity of the classical versus modified algorithms are the ideal companion to
\cref{sec:cgsmgs,sec:lossorthog}. For the full algebraic and floating-point analysis---
the $\kappa^2 u$ versus $\kappa u$ loss-of-orthogonality bounds, the ``twice is
enough'' re-orthogonalization result, and the Arnoldi and Lanczos recurrences with
their reorthogonalization variants---Golub and Van Loan's \emph{Matrix
Computations}~\citep{golubvanloan2013} is the standard reference. For the eigenvalue
side specifically---the ghost-eigenvalue phenomenon, selective and partial
reorthogonalization, and the practical Lanczos and Arnoldi eigensolvers that
\cref{ch:krylovschur} builds on---Saad's \emph{Numerical Methods for Large Eigenvalue
Problems}~\citep{saad2011eig} is the place to go deeper. Read all three with the
\emph{orthogonality contract} of \cref{sec:contract} in mind: each is, at bottom, a
study of how well that contract can be held and at what cost.

\exercisessection
\begin{enumerate}[nosep]
  \item \emph{(Geometry.)} In $\Real^3$, let $\vq_1 = (1,0,0)^{\!\top}$,
  $\vq_2 = (0,1,0)^{\!\top}$, and $\vw = (3, 4, 5)^{\!\top}$. Compute the
  Gram--Schmidt residual $\vw' = \vw - \inner{\vw}{\vq_1}\vq_1 -
  \inner{\vw}{\vq_2}\vq_2$ and verify by direct inner product that
  $\inner{\vw'}{\vq_1} = \inner{\vw'}{\vq_2} = 0$ (the contract of
  \cref{lem:residorthog}). What is $\vq_3$ after normalization?
  \item \emph{(CGS versus MGS.)} Repeat \cref{wex:cgsmgs} with $\varepsilon = 10^{-3}$
  and confirm that the residual orthogonality is $\sim\varepsilon$ for CGS and
  $\sim\varepsilon^2$ for MGS. Then explain in one sentence why MGS's second
  projection coefficient $H_2$ was nonzero while CGS's was (numerically) zero, in
  terms of \emph{which vector each coefficient was taken against}.
  \item \emph{(The sequential obstruction.)} Explain why classical Gram--Schmidt can
  batch its $m$ inner products into a single matrix-vector product
  $\mat{Q}^{\!\top}\vw$ but modified Gram--Schmidt cannot. Which combinator
  (\cref{ch:combinators}) is each---\code{map} or \code{fold}? Why does permuting the
  columns of $\mat{Q}$ change the MGS result (at the bit level) but not the CGS result?
  \item \emph{(Twice is enough.)} A first Gram--Schmidt pass leaves a residual with a
  relative overlap of $\sim u\kappa$ against the span. Argue heuristically why a second
  pass reduces the overlap to $\sim u\cdot(u\kappa)$ and hence why a third pass is
  unnecessary. Under what condition on $\kappa$ does even \emph{one} pass already
  suffice (so that CGS2 is overkill)?
  \item \emph{(Arnoldi shape.)} The Arnoldi relation is
  $\mat{A}\mat{V}_m = \mat{V}_{m+1}\bar{\mat{H}}_m$. Give the dimensions of each of the
  three matrices, and explain why $\bar{\mat{H}}_m$ has exactly one more row than
  column and is zero below its first sub-diagonal. Which entry of $\bar{\mat{H}}_m$ is
  the residual norm $\norm{\vw'}$ from the normalization step, and which chapter
  (\cref{ch:gmres} or \cref{ch:krylovschur}) reads which feature of this matrix?
  \item \emph{(Lanczos by hand.)} Run two Lanczos steps as in \cref{wex:lanczos} on
  $\mat{A} = \begin{psmallmatrix} 4 & 1 & 0 \\ 1 & 4 & 1 \\ 0 & 1 & 4 \end{psmallmatrix}$
  with $\vq_1 = (1,0,0)^{\!\top}$. Compute $\alpha_1, \beta_1, \vq_2, \alpha_2,
  \beta_2, \vq_3$, assemble $\mat{T}_2$, and verify that $\vq_3$ is automatically
  orthogonal to $\vq_1$ even though you never projected against $\vq_1$ in step~2.
  \item \emph{(Why symmetry shortens the recurrence.)} Prove that the Arnoldi
  Hessenberg matrix is symmetric when $\mat{A}$ is symmetric, and conclude that it is
  tridiagonal. Then state precisely which orthogonalization coefficients Lanczos is
  entitled to skip, and explain---in terms of finite-precision arithmetic---why
  skipping them is exactly what causes Lanczos to lose orthogonality and produce ghost
  eigenvalues.
  \item \emph{(The contract as an interface.)} \cref{ch:rotations} calls
  \lstinline[style=l4style]@orthog := mgs | cgs | cgs2@ a genuine rotation by the
  algorithmic-substitution test. Identify the contract that all three satisfy and
  argue that a GMRES solver written against that contract is correct regardless of
  which variant is installed. Give one concrete way the \emph{choice} of variant could
  still affect the solver's observable behavior, even though all three satisfy the
  contract. (Hint: \cref{sec:lossorthog}.)
\end{enumerate}
